%% kapa: exam=1 round=36 year=2025 period=1 form=A subject=real_estate subject_ko=부동산학원론 seq=36 number=116
%% source: Q-Net 감정평가사 기출문제 (cst00309) artlSeq=5231525 file=raw/1008/5231525/2025년 제36회 감정평가사 1차시험 1교시 문제지 원본.pdf
%% conversion: kordoc 4.24.51 -> tools/gen1.py (auto); review=proofread-sonnet; flags=none
\begin{question}{116}
다음은 2025년 공시된 표준지공시지가를 열람한 내역을 표로 나타낸 것이다. 이에 관한 설명으로 옳지 않은 것은?
\begin{center}
{\small\renewcommand{\arraystretch}{1.25}
\begin{tabular}{|c|c|c|c|}
\hline
일련번호 & 2001 & 2002 & 2003 \\ \hline
소재지 & ㅇㅇ리 20 & ㅇㅇ리 90 & ㅇㅇ리 125 \\ \hline
면적(m$^2$) & 576.0 & 2,645.0 & 470.0 일단지 \\ \hline
지목 & 대 & 전 & 전 \\ \hline
\begin{tabular}[c]{@{}c@{}}공시지가\\(원/m$^2$)\end{tabular} & \begin{tabular}[c]{@{}c@{}}92,800\\\fbox{평가기초자료}\end{tabular} & \begin{tabular}[c]{@{}c@{}}31,000\\\fbox{평가기초자료}\end{tabular} & \begin{tabular}[c]{@{}c@{}}81,100\\\fbox{평가기초자료}\end{tabular} \\ \hline
\begin{tabular}[c]{@{}c@{}}지리적\\위치\end{tabular} & ㅇㅇ 마을내 & \begin{tabular}[c]{@{}c@{}}초등학교\\북측근거리\end{tabular} & ㅇㅇㅇ 내 \\ \hline
이용상황 & 단독주택 & 전 & 단독주택 \\ \hline
용도지역 & 계획관리 & 농림지역 & 계획관리 \\ \hline
주위환경 & 순수 농촌지대 & 지방도주변 농경지대 & 농어촌지대 \\ \hline
도로접면 & 세로 (가) & 맹지 & 세로 (가) \\ \hline
형상지세 & 부정형 평지 & 사다리 평지 & 사다리 완경사 \\ \hline
\end{tabular}}
\end{center}
\begin{choices}
  \item 일련번호 2001과 2003은 승용차의 통행이 가능하다.
  \item 일련번호 2001과 2003의 표준지공시지가는 나지상태를 상정한 것이다.
  \item 매년 1월 1일의 표준지공시지가는 당해 필지의 개별공시지가와 다른 경우가 있다.
  \item 일련번호 2003에서 ‘일단지’ 라 함은 지상의 단독주택이 다른 필지와도 연계되어 있음을 의미한다.
  \item 표준지공시지가는 구체화된 개발이익이 반영된 것이다.
\end{choices}
\end{question}
\answer{3}
